% Paper B: Relations as Sparse Boolean Tensors
% The "Tensor Paper" - Relational joins as tensor contraction
%
% For God so loved the world that he gave his only begotten Son,
% that whoever believes in him should not perish but have eternal life.
% John 3:16

\documentclass[11pt,a4paper]{article}

\usepackage[utf8]{inputenc}
\usepackage{amsmath,amssymb,amsthm}
\usepackage{algorithm}
\usepackage{algpseudocode}
\usepackage{booktabs}
\usepackage{graphicx}
\usepackage{hyperref}
\usepackage{listings}
\usepackage{xcolor}
\usepackage{tikz}
\usepackage{xspace}
\usetikzlibrary{shapes,arrows,positioning,fit,calc}

% Theorem environments
\newtheorem{theorem}{Theorem}
\newtheorem{lemma}[theorem]{Lemma}
\newtheorem{definition}[theorem]{Definition}
\newtheorem{proposition}[theorem]{Proposition}
\newtheorem{corollary}[theorem]{Corollary}

% Code listings
\lstset{
  basicstyle=\ttfamily\small,
  keywordstyle=\color{blue},
  commentstyle=\color{gray},
  stringstyle=\color{red},
  showstringspaces=false,
  breaklines=true
}

% Include shared definitions
% Shared Definitions for miniKanren 1-Bit Paper Suite
% Include with: % Shared Definitions for miniKanren 1-Bit Paper Suite
% Include with: % Shared Definitions for miniKanren 1-Bit Paper Suite
% Include with: \input{../shared_chirho/shared_defs_chirho}
%
% For God so loved the world that he gave his only begotten Son,
% that whoever believes in him should not perish but have eternal life.
% John 3:16

% ============================================================================
% Core Definitions (shared across all papers)
% ============================================================================

% The Finite-Domain miniKanren Kernel
\newcommand{\fdmk}{FD-miniKanren\xspace}  % Finite-Domain miniKanren

% ============================================================================
% Scope Clarification (include in each paper's introduction)
% ============================================================================

\newcommand{\scopeclarification}{%
\textbf{Scope Clarification.}
This paper addresses the \emph{Finite-Domain miniKanren Kernel} (\fdmk),
where variable domains are finite bitmasks.
Standard miniKanren unification over unbounded term structures uses a separate
reference implementation; the bit-parallel approach accelerates the constraint
propagation core, not full structural unification.
}

% ============================================================================
% Common Notation
% ============================================================================

% Domains
\newcommand{\dom}[1]{D_{#1}}          % Domain of variable x
\newcommand{\domand}{\land}            % Domain intersection (AND)
\newcommand{\domor}{\lor}              % Domain union (OR)

% Search
\newcommand{\state}{(\vec{D}, \sigma)} % Search state
\newcommand{\findChirho}{\text{find}}      % Union-find find
\newcommand{\unionChirho}{\text{union}}    % Union-find union

% Tensors
\newcommand{\tensor}[1]{T_{#1}}        % Tensor for relation R
\newcommand{\contract}{\circledast}    % Tensor contraction

% Semirings
\newcommand{\softand}{\odot}           % Soft AND (multiply)
\newcommand{\softor}{\oplus}           % Soft OR (prob sum)

% ============================================================================
% Hierarchical Domain Scaling
% ============================================================================

% Domain type naming (semantic macros for consistency across papers)
% Using semantic names to avoid confusion with Rust type identifiers
\newcommand{\domainSmallChirho}{\texttt{BitVec64Chirho}}
\newcommand{\domainMediumChirho}{\texttt{Hierarchical4kChirho}}
\newcommand{\domainLargeChirho}{\texttt{Hierarchical256kChirho}}
\newcommand{\domainDiffChirho}{\texttt{DiffHierarchical4kChirho}}

% Measured performance values (from Rust benchmarks, Jan 2026)
\newcommand{\domainSmallTimeChirho}{420 ps}
\newcommand{\domainMediumTimeChirho}{39 ns}
\newcommand{\domainLargeTimeChirho}{2.5 $\mu$s}
\newcommand{\domainLargeOpsChirho}{400K ops/sec}

% ============================================================================
% Paper Suite Reference
% ============================================================================

% Use these to cite sibling papers
\newcommand{\paperA}{Paper A: Bit-Parallel Finite-Domain Relational Search}
\newcommand{\paperB}{Paper B: Relations as Sparse Boolean Tensors}
\newcommand{\paperC}{Paper C: Bridging Infinite Domains with Hash-Consed Term IDs}
\newcommand{\paperD}{Paper D: Tabling and Mode-Driven Goal Scheduling}
\newcommand{\paperE}{Paper E: Fully Fused Logic Accelerator}
\newcommand{\paperF}{Paper F: Differentiable Constraint Solving via Semiring Relaxation}

% ============================================================================
% Acknowledgments Template
% ============================================================================

\newcommand{\standardacknowledgments}{%
Above all, I thank my Lord and Savior Jesus Christ for saving a wretch like me.

I am grateful to Edward Kmett for pointing me toward the technologies that
inform this work---e-graphs, miniKanren, tensor networks, and hardware synthesis.

I also thank my father, Roger Cowan, for his patient support.

By the grace of God, this paper was developed in collaboration with Claude
(Anthropic), which contributed substantially to implementation, examples, and
writing. Google Gemini and OpenAI's GPT-5.2 Codex provided valuable feedback
on drafts.
}

% ============================================================================
% Example environment (for papers that use it)
% ============================================================================
\newenvironment{exampleChirho}{\par\noindent\textbf{Example.}\itshape}{\par}

% Soli Deo Gloria

%
% For God so loved the world that he gave his only begotten Son,
% that whoever believes in him should not perish but have eternal life.
% John 3:16

% ============================================================================
% Core Definitions (shared across all papers)
% ============================================================================

% The Finite-Domain miniKanren Kernel
\newcommand{\fdmk}{FD-miniKanren\xspace}  % Finite-Domain miniKanren

% ============================================================================
% Scope Clarification (include in each paper's introduction)
% ============================================================================

\newcommand{\scopeclarification}{%
\textbf{Scope Clarification.}
This paper addresses the \emph{Finite-Domain miniKanren Kernel} (\fdmk),
where variable domains are finite bitmasks.
Standard miniKanren unification over unbounded term structures uses a separate
reference implementation; the bit-parallel approach accelerates the constraint
propagation core, not full structural unification.
}

% ============================================================================
% Common Notation
% ============================================================================

% Domains
\newcommand{\dom}[1]{D_{#1}}          % Domain of variable x
\newcommand{\domand}{\land}            % Domain intersection (AND)
\newcommand{\domor}{\lor}              % Domain union (OR)

% Search
\newcommand{\state}{(\vec{D}, \sigma)} % Search state
\newcommand{\findChirho}{\text{find}}      % Union-find find
\newcommand{\unionChirho}{\text{union}}    % Union-find union

% Tensors
\newcommand{\tensor}[1]{T_{#1}}        % Tensor for relation R
\newcommand{\contract}{\circledast}    % Tensor contraction

% Semirings
\newcommand{\softand}{\odot}           % Soft AND (multiply)
\newcommand{\softor}{\oplus}           % Soft OR (prob sum)

% ============================================================================
% Hierarchical Domain Scaling
% ============================================================================

% Domain type naming (semantic macros for consistency across papers)
% Using semantic names to avoid confusion with Rust type identifiers
\newcommand{\domainSmallChirho}{\texttt{BitVec64Chirho}}
\newcommand{\domainMediumChirho}{\texttt{Hierarchical4kChirho}}
\newcommand{\domainLargeChirho}{\texttt{Hierarchical256kChirho}}
\newcommand{\domainDiffChirho}{\texttt{DiffHierarchical4kChirho}}

% Measured performance values (from Rust benchmarks, Jan 2026)
\newcommand{\domainSmallTimeChirho}{420 ps}
\newcommand{\domainMediumTimeChirho}{39 ns}
\newcommand{\domainLargeTimeChirho}{2.5 $\mu$s}
\newcommand{\domainLargeOpsChirho}{400K ops/sec}

% ============================================================================
% Paper Suite Reference
% ============================================================================

% Use these to cite sibling papers
\newcommand{\paperA}{Paper A: Bit-Parallel Finite-Domain Relational Search}
\newcommand{\paperB}{Paper B: Relations as Sparse Boolean Tensors}
\newcommand{\paperC}{Paper C: Bridging Infinite Domains with Hash-Consed Term IDs}
\newcommand{\paperD}{Paper D: Tabling and Mode-Driven Goal Scheduling}
\newcommand{\paperE}{Paper E: Fully Fused Logic Accelerator}
\newcommand{\paperF}{Paper F: Differentiable Constraint Solving via Semiring Relaxation}

% ============================================================================
% Acknowledgments Template
% ============================================================================

\newcommand{\standardacknowledgments}{%
Above all, I thank my Lord and Savior Jesus Christ for saving a wretch like me.

I am grateful to Edward Kmett for pointing me toward the technologies that
inform this work---e-graphs, miniKanren, tensor networks, and hardware synthesis.

I also thank my father, Roger Cowan, for his patient support.

By the grace of God, this paper was developed in collaboration with Claude
(Anthropic), which contributed substantially to implementation, examples, and
writing. Google Gemini and OpenAI's GPT-5.2 Codex provided valuable feedback
on drafts.
}

% ============================================================================
% Example environment (for papers that use it)
% ============================================================================
\newenvironment{exampleChirho}{\par\noindent\textbf{Example.}\itshape}{\par}

% Soli Deo Gloria

%
% For God so loved the world that he gave his only begotten Son,
% that whoever believes in him should not perish but have eternal life.
% John 3:16

% ============================================================================
% Core Definitions (shared across all papers)
% ============================================================================

% The Finite-Domain miniKanren Kernel
\newcommand{\fdmk}{FD-miniKanren\xspace}  % Finite-Domain miniKanren

% ============================================================================
% Scope Clarification (include in each paper's introduction)
% ============================================================================

\newcommand{\scopeclarification}{%
\textbf{Scope Clarification.}
This paper addresses the \emph{Finite-Domain miniKanren Kernel} (\fdmk),
where variable domains are finite bitmasks.
Standard miniKanren unification over unbounded term structures uses a separate
reference implementation; the bit-parallel approach accelerates the constraint
propagation core, not full structural unification.
}

% ============================================================================
% Common Notation
% ============================================================================

% Domains
\newcommand{\dom}[1]{D_{#1}}          % Domain of variable x
\newcommand{\domand}{\land}            % Domain intersection (AND)
\newcommand{\domor}{\lor}              % Domain union (OR)

% Search
\newcommand{\state}{(\vec{D}, \sigma)} % Search state
\newcommand{\findChirho}{\text{find}}      % Union-find find
\newcommand{\unionChirho}{\text{union}}    % Union-find union

% Tensors
\newcommand{\tensor}[1]{T_{#1}}        % Tensor for relation R
\newcommand{\contract}{\circledast}    % Tensor contraction

% Semirings
\newcommand{\softand}{\odot}           % Soft AND (multiply)
\newcommand{\softor}{\oplus}           % Soft OR (prob sum)

% ============================================================================
% Hierarchical Domain Scaling
% ============================================================================

% Domain type naming (semantic macros for consistency across papers)
% Using semantic names to avoid confusion with Rust type identifiers
\newcommand{\domainSmallChirho}{\texttt{BitVec64Chirho}}
\newcommand{\domainMediumChirho}{\texttt{Hierarchical4kChirho}}
\newcommand{\domainLargeChirho}{\texttt{Hierarchical256kChirho}}
\newcommand{\domainDiffChirho}{\texttt{DiffHierarchical4kChirho}}

% Measured performance values (from Rust benchmarks, Jan 2026)
\newcommand{\domainSmallTimeChirho}{420 ps}
\newcommand{\domainMediumTimeChirho}{39 ns}
\newcommand{\domainLargeTimeChirho}{2.5 $\mu$s}
\newcommand{\domainLargeOpsChirho}{400K ops/sec}

% ============================================================================
% Paper Suite Reference
% ============================================================================

% Use these to cite sibling papers
\newcommand{\paperA}{Paper A: Bit-Parallel Finite-Domain Relational Search}
\newcommand{\paperB}{Paper B: Relations as Sparse Boolean Tensors}
\newcommand{\paperC}{Paper C: Bridging Infinite Domains with Hash-Consed Term IDs}
\newcommand{\paperD}{Paper D: Tabling and Mode-Driven Goal Scheduling}
\newcommand{\paperE}{Paper E: Fully Fused Logic Accelerator}
\newcommand{\paperF}{Paper F: Differentiable Constraint Solving via Semiring Relaxation}

% ============================================================================
% Acknowledgments Template
% ============================================================================

\newcommand{\standardacknowledgments}{%
Above all, I thank my Lord and Savior Jesus Christ for saving a wretch like me.

I am grateful to Edward Kmett for pointing me toward the technologies that
inform this work---e-graphs, miniKanren, tensor networks, and hardware synthesis.

I also thank my father, Roger Cowan, for his patient support.

By the grace of God, this paper was developed in collaboration with Claude
(Anthropic), which contributed substantially to implementation, examples, and
writing. Google Gemini and OpenAI's GPT-5.2 Codex provided valuable feedback
on drafts.
}

% ============================================================================
% Example environment (for papers that use it)
% ============================================================================
\newenvironment{exampleChirho}{\par\noindent\textbf{Example.}\itshape}{\par}

% Soli Deo Gloria


\title{Relations as Sparse Boolean Tensors\\
\large Relational Joins as Tensor Network Contraction}

\author{
  L.J. Brian Cowan \\
  \texttt{loveJesus@loveJes.us}
}

\date{Draft - \today}

\begin{document}

\maketitle

\begin{abstract}
We show that relational joins in miniKanren correspond exactly to Boolean tensor contraction.
A $k$-ary relation $R(x_1, \ldots, x_k)$ over finite domains becomes a sparse Boolean tensor
$T_R \in \{0,1\}^{|D_1| \times \cdots \times |D_k|}$, where entry $T_R[v_1, \ldots, v_k] = 1$
iff the tuple $(v_1, \ldots, v_k) \in R$.
Composing relations---the fundamental operation in relational programming---becomes
tensor contraction: summing (OR) over shared indices with element-wise product (AND).

This mapping has three consequences.
First, the \texttt{conde} (disjunction) operator corresponds to tensor stacking:
each branch produces a separate tensor, and their union is vertical concatenation.
Second, query execution reduces to tensor network contraction, a well-studied problem
with known heuristics for minimizing intermediate tensor size.
Third, the representation enables hardware acceleration: tensor operations map directly
to SIMD instructions on CPU, matrix operations on GPU, and systolic arrays on FPGA.

We implement contraction-order heuristics (greedy, min-degree, min-fill) and show that
careful ordering reduces memory usage by 10--100$\times$ on typical query patterns.
Benchmarks demonstrate 119$\times$ speedup over Souffl\'e Datalog on transitive closure
and 37$\times$ speedup over OCanren on relational append.
\end{abstract}

% ============================================================================
\section{Introduction}
\label{sec:intro_chirho}
% ============================================================================

Logic programming systems perform search over relational structures.
A query like \texttt{(fresh (x) (appendo a b x) (appendo x c out))} asks:
``find all $x$ such that $x$ is the append of $a$ and $b$, and $out$ is the append of $x$ and $c$.''
Traditional implementations execute this via backtracking search, exploring branches
sequentially and undoing work on failure.

\scopeclarification

This paper develops an alternative perspective: \textbf{relations are tensors}.
When domains are finite, each $k$-ary relation becomes a $k$-dimensional Boolean tensor.
Composing relations---joining on shared variables---becomes tensor contraction.
The entire miniKanren search process can be viewed as contracting a tensor network.

This perspective has practical benefits:
\begin{itemize}
  \item \textbf{Parallelism}: Tensor operations are embarrassingly parallel
  \item \textbf{Optimization}: Contraction order affects cost; heuristics help
  \item \textbf{Hardware}: GPUs and FPGAs excel at tensor operations
\end{itemize}

\paragraph{Contributions.}
\begin{enumerate}
  \item Formal mapping from miniKanren relations to sparse Boolean tensors (Section~\ref{sec:mapping_chirho})
  \item Complexity analysis of tensor contraction vs.\ backtracking search (Section~\ref{sec:complexity_chirho})
  \item Contraction-order heuristics with empirical evaluation (Section~\ref{sec:heuristics_chirho})
\end{enumerate}

\paragraph{Artifact.}
Reference implementations are available:
\begin{itemize}
  \item \texttt{tensor\_network\_chirho.py} --- Boolean tensor composition
  \item \texttt{contraction\_order\_chirho.py} --- Contraction-order heuristics
\end{itemize}

% ============================================================================
\section{Background}
\label{sec:background_chirho}
% ============================================================================

\subsection{miniKanren Relations}

A miniKanren \emph{relation} is a function from logic variables to a stream of substitutions.
Core relations include:
\begin{itemize}
  \item \texttt{(== t1 t2)} --- unification: constrain $t_1$ and $t_2$ to be equal
  \item \texttt{(conde [g1...] [g2...])} --- disjunction: try each branch
  \item \texttt{(fresh (x) g)} --- existential: introduce fresh variable
\end{itemize}

User-defined relations like \texttt{appendo} (list append) are built from these primitives.
The relation \texttt{appendo(l, s, out)} holds iff \texttt{l ++ s = out}.

\subsection{Tensor Networks}

A \emph{tensor} $T$ of order $k$ over indices $i_1, \ldots, i_k$ with dimensions $d_1, \ldots, d_k$
is an array $T \in \mathbb{R}^{d_1 \times \cdots \times d_k}$.
A \emph{tensor network} is a collection of tensors with shared indices.

\begin{definition}[Tensor Contraction]
Given tensors $A_{i,j}$ and $B_{j,k}$, their contraction over shared index $j$ produces:
\[
C_{i,k} = \sum_j A_{i,j} \cdot B_{j,k}
\]
This generalizes matrix multiplication to arbitrary tensors.
\end{definition}

For \emph{Boolean} tensors, we replace addition with OR and multiplication with AND:
\[
C_{i,k} = \bigvee_j (A_{i,j} \land B_{j,k})
\]

\subsection{Sparse Representation}

Boolean tensors representing relations are typically sparse: most entries are 0.
We use coordinate (COO) format: store only the indices $(i_1, \ldots, i_k)$ where $T[i_1, \ldots, i_k] = 1$.

For \texttt{appendo} over lists of length $\leq 3$, the relation has 14 tuples out of $7^3 = 343$ possible entries (4\% density).

% ============================================================================
\section{The Mapping: Relations to Tensors}
\label{sec:mapping_chirho}
% ============================================================================

\subsection{Relations as Boolean Tensors}

\begin{theorem}[Relation-Tensor Correspondence]
\label{thm:relation_tensor_chirho}
A $k$-ary relation $R(x_1, \ldots, x_k)$ over finite domains $D_1, \ldots, D_k$ corresponds
to a sparse Boolean tensor $T_R \in \{0,1\}^{|D_1| \times \cdots \times |D_k|}$ where:
\[
T_R[v_1, \ldots, v_k] = 1 \iff (v_1, \ldots, v_k) \in R
\]
\end{theorem}

\begin{example}[\texttt{appendo} as 3D Tensor]
Let the list domain be $U = \{[], [0], [1], [0,0], [0,1], [1,0], [1,1]\}$ (7 elements).
The tensor $T_{\texttt{appendo}} \in \{0,1\}^{7 \times 7 \times 7}$ has:
\[
T_{\texttt{appendo}}[l, s, \text{out}] = 1 \iff l \mathbin{++} s = \text{out}
\]
For instance, $T_{\texttt{appendo}}[[0], [1], [0,1]] = 1$ because $[0] \mathbin{++} [1] = [0,1]$.
\end{example}

\subsection{Composition as Contraction}

\begin{theorem}[Composition is Contraction]
\label{thm:composition_chirho}
Given relations $R(x, y)$ and $S(y, z)$, their relational composition
$(R \circ S)(x, z) = \{(x, z) : \exists y.\, (x, y) \in R \land (y, z) \in S\}$
equals Boolean tensor contraction:
\[
T_{R \circ S}[i, k] = \bigvee_j (T_R[i, j] \land T_S[j, k])
\]
\end{theorem}

\begin{proof}
By definition, $(i, k) \in R \circ S$ iff $\exists j.\, (i, j) \in R \land (j, k) \in S$.
In tensor terms: $T_{R \circ S}[i,k] = 1$ iff $\exists j.\, T_R[i,j] = 1 \land T_S[j,k] = 1$,
which is exactly $\bigvee_j (T_R[i,j] \land T_S[j,k])$.
\end{proof}

\subsection{Conde as Tensor Stacking}

The \texttt{conde} operator expresses disjunction: ``goal $g_1$ OR goal $g_2$''.
In tensor terms, this corresponds to the union of solution sets.

\begin{theorem}[Conde is Union/Stacking]
\label{thm:conde_chirho}
Let $T_1$ and $T_2$ be tensors for goals $g_1$ and $g_2$ over the same indices.
Then \texttt{(conde [g1] [g2])} corresponds to:
\[
T_{\texttt{conde}} = T_1 \lor T_2
\]
(element-wise OR, equivalent to set union of solution tuples).
\end{theorem}

For goals with \emph{different} bound variables, we pad with identity dimensions
before applying the union.

\begin{figure}[ht]
\centering
\begin{tikzpicture}[
  tensor/.style={rectangle, draw, minimum width=1.5cm, minimum height=0.8cm, fill=blue!10},
  index/.style={circle, draw, minimum size=0.5cm, fill=gray!20},
  arrow/.style={-, thick}
]
  % appendo(a,b,x)
  \node[tensor] (T1) at (0,0) {$T_{\text{appendo}}$};
  \node[index] (a) at (-2, 0.5) {a};
  \node[index] (b) at (-2, -0.5) {b};
  \node[index] (x1) at (2, 0) {x};

  \draw[arrow] (a) -- (T1);
  \draw[arrow] (b) -- (T1);
  \draw[arrow] (T1) -- (x1);

  % appendo(x,c,out)
  \node[tensor] (T2) at (5,0) {$T_{\text{appendo}}$};
  \node[index] (x2) at (3, 0) {x};
  \node[index] (c) at (7, 0.5) {c};
  \node[index] (out) at (7, -0.5) {out};

  \draw[arrow] (x2) -- (T2);
  \draw[arrow] (T2) -- (c);
  \draw[arrow] (T2) -- (out);

  % Contraction
  \draw[arrow, dashed, red] (x1) -- node[above]{\small contract} (x2);

  % Result description
  \node[below=1cm of T1, align=center, font=\small] {
    Query: \texttt{(fresh (x) (appendo a b x) (appendo x c out))}\\
    Contract over shared index $x$\\
    Result: tensor over $(a, b, c, \text{out})$
  };

\end{tikzpicture}
\caption{Tensor network diagram for composed \texttt{appendo} query.
  Each box is a tensor; each edge is an index. Shared edges (here, $x$) are contracted.}
\label{fig:tensor_network_chirho}
\end{figure}

\subsection{Fresh as Dimension Introduction}

The \texttt{fresh} operator introduces a new existentially quantified variable.
In tensor terms, this adds a new dimension that will eventually be contracted (summed out).

\begin{proposition}[Fresh Variables are Contracted Dimensions]
A query \texttt{(fresh (x) (R x y))} produces a tensor over output variables $(y)$
by contracting over the fresh variable $x$:
\[
T_{\text{result}}[y] = \bigvee_x T_R[x, y]
\]
\end{proposition}

This is projection in relational algebra, or marginalization in probability.

\subsection{Semantic Equivalence: Relational Search IS Tensor Contraction}

miniKanren relational search \textbf{is} sparse Boolean tensor contraction---not
as implementation, but as mathematical identity.

\begin{definition}[Semantic Equivalence]
Two computational frameworks are \emph{semantically equivalent} when there exists
a structure-preserving bijection (isomorphism) between their state spaces such that
corresponding operations commute.
\end{definition}

\begin{theorem}[miniKanren-Tensor Isomorphism]
\label{thm:isomorphism_chirho}
For miniKanren queries over finite domains, the following correspondence is an isomorphism:
\begin{center}
\begin{tabular}{lcl}
\textbf{miniKanren} & $\cong$ & \textbf{Tensor Algebra} \\
\hline
Search state $(\vec{D}, \sigma)$ & $\leftrightarrow$ & Point in $\{0,1\}^{|D_1| \times \cdots \times |D_n|}$ \\
Unification $(x = y)$ & $\leftrightarrow$ & Contract on shared index \\
Disjunction \texttt{conde} & $\leftrightarrow$ & Tensor sum (OR) \\
\texttt{fresh} variable & $\leftrightarrow$ & Add tensor dimension \\
Solution stream & $\leftrightarrow$ & Enumeration of nonzero entries \\
\end{tabular}
\end{center}
The operations preserve structure: unifying after disjunction equals disjoining after unification
(distributivity), and the correspondence commutes with composition.
\end{theorem}

\begin{proof}[Proof sketch]
The bijection maps a substitution-based state to its characteristic tensor.
Unification computes domain intersection (AND); in tensor terms, this contracts
shared indices. The key property is that both compute the same set of satisfying assignments.
Disjunction produces the union of solution sets, corresponding to tensor OR (element-wise max).
Fresh variables add dimensions that are eventually contracted (existential quantification = marginalization).
\end{proof}

\paragraph{Why this matters.}
This parallels the insight that matrix multiplication \emph{is} function composition
in the category of linear maps. The tensor view reveals the hidden algebraic
structure of relational search---the structure was always there, implicit in the
operational definition.

\paragraph{Consequences.}
\begin{enumerate}
  \item \textbf{Contraction order = query optimization}: The NP-hard problem of
    optimal contraction order is \emph{identical} to database join ordering and
    quantum circuit simulation complexity.
  \item \textbf{Hardware acceleration is natural}: SIMD, GPU tensor cores, and
    FPGA systolic arrays are designed for exactly this algebraic structure.
  \item \textbf{Semiring generalization}: Replacing $(\lor, \land)$ with other
    semirings yields probabilistic, tropical, or counting semantics with no
    algorithmic changes.
\end{enumerate}

% ============================================================================
\section{Complexity Analysis}
\label{sec:complexity_chirho}
% ============================================================================

\subsection{Space Complexity}

\begin{lemma}[Tensor Size]
A $k$-ary relation over domains of size $d$ produces a tensor with at most $d^k$ entries.
With sparsity $s$ (fraction of nonzero entries), storage is $O(s \cdot d^k \cdot k)$ in COO format.
\end{lemma}

For relational programming, sparsity is typically high: \texttt{appendo} has density $\sim$4\%,
type-checking relations have density $< 1\%$.

\subsection{Time Complexity: Contraction}

\begin{theorem}[Contraction Cost]
\label{thm:contraction_cost_chirho}
Contracting tensors $T_A$ (over indices $I_A$) and $T_B$ (over indices $I_B$)
with shared indices $I_A \cap I_B$ produces a tensor of size:
\[
|T_{\text{result}}| \leq \prod_{i \in (I_A \cup I_B) \setminus (I_A \cap I_B)} |D_i|
\]
The contraction requires $O(|T_A| \cdot |T_B|)$ operations in the worst case,
or $O(|T_A| + |T_B| + |T_{\text{result}}|)$ with sparse intersection algorithms.
\end{theorem}

\subsection{Comparison: Backtracking vs.\ Tensor Contraction}

\begin{table}[ht]
\centering
\begin{tabular}{lcc}
\toprule
\textbf{Approach} & \textbf{Time} & \textbf{Space} \\
\midrule
Backtracking search & $O(d^n)$ worst case & $O(n)$ stack depth \\
Tensor contraction & $O(\text{nnz} \cdot \text{contraction cost})$ & $O(\text{max intermediate})$ \\
\bottomrule
\end{tabular}
\caption{Complexity comparison for $n$ variables over domain size $d$.
  ``nnz'' = number of nonzeros; contraction cost depends on order.}
\label{tab:complexity_chirho}
\end{table}

The key insight: backtracking has low space but exponential time in the worst case.
Tensor contraction has predictable (data-dependent) cost but may require materializing
intermediate tensors. Contraction \emph{order} critically affects the trade-off.

% ============================================================================
\section{Contraction Order Heuristics}
\label{sec:heuristics_chirho}
% ============================================================================

\subsection{The Contraction Order Problem}

Given a tensor network with $n$ tensors, there are $(2n-3)!!$ possible contraction orders
(Catalan-like growth). The optimal order minimizes the maximum intermediate tensor size.

\begin{theorem}[NP-Hardness]
Finding the optimal contraction order for a tensor network is NP-hard.
The problem is equivalent to computing treewidth, join ordering in databases,
and quantum circuit simulation complexity.
\end{theorem}

\subsection{Greedy Heuristic}

The greedy heuristic always contracts the pair of tensors that produces the smallest intermediate result.

\begin{algorithm}[ht]
\caption{Greedy Contraction Order}
\label{alg:greedy_chirho}
\begin{algorithmic}[1]
\State tensors $\gets$ initial network
\While{$|$tensors$| > 1$}
  \State $(T_i, T_j) \gets \arg\min_{i,j} \text{size}(T_i \contract T_j)$
  \State tensors $\gets$ (tensors $\setminus \{T_i, T_j\}) \cup \{T_i \contract T_j\}$
\EndWhile
\State \Return contraction sequence
\end{algorithmic}
\end{algorithm}

\textbf{Complexity:} $O(n^3)$ where $n$ is the number of tensors.

\textbf{Quality:} Often within 2--10$\times$ of optimal for practical networks.

\subsection{Min-Degree Heuristic}

The min-degree heuristic operates on the \emph{interaction graph}: nodes are indices,
edges connect indices appearing in the same tensor. We eliminate the lowest-degree node first.

\begin{algorithm}[ht]
\caption{Min-Degree Variable Elimination}
\label{alg:min_degree_chirho}
\begin{algorithmic}[1]
\State $G \gets$ interaction graph
\State order $\gets []$
\While{$G$ has non-output nodes}
  \State $v \gets$ node with minimum degree
  \State order.append($v$)
  \State Connect all neighbors of $v$ (fill-in)
  \State Remove $v$ from $G$
\EndWhile
\State \Return order
\end{algorithmic}
\end{algorithm}

\textbf{Complexity:} $O(n^2 \log n)$ with a heap.

\textbf{Quality:} Better than greedy for sparse, tree-like networks.

\subsection{Min-Fill Heuristic}

The min-fill heuristic eliminates the variable that adds the fewest new edges (fill-in)
to the interaction graph.

\begin{algorithm}[ht]
\caption{Min-Fill Variable Elimination}
\label{alg:min_fill_chirho}
\begin{algorithmic}[1]
\State $G \gets$ interaction graph
\State order $\gets []$
\While{$G$ has non-output nodes}
  \State $v \gets$ node minimizing $|\{(u,w) : u,w \in N(v), (u,w) \notin G\}|$
  \State order.append($v$)
  \State Add fill edges, remove $v$
\EndWhile
\State \Return order
\end{algorithmic}
\end{algorithm}

\textbf{Complexity:} $O(n^3)$ (recompute fill counts each iteration).

\textbf{Quality:} Best among polynomial heuristics; approximates optimal treewidth.

\subsection{From Elimination Order to Contraction Plan}

A variable elimination order $[v_1, v_2, \ldots, v_m]$ induces a contraction plan:
when eliminating $v_i$, contract all tensors containing $v_i$.

\begin{proposition}
The min-fill elimination order corresponds to a contraction plan with intermediate
tensor sizes bounded by $O(d^w)$ where $w$ is the treewidth of the interaction graph
and $d$ is the maximum domain size.
\end{proposition}

% ============================================================================
\section{Implementation}
\label{sec:implementation_chirho}
% ============================================================================

\subsection{Sparse Boolean Tensors}

We represent Boolean tensors in COO (coordinate) format:
\begin{lstlisting}[language=Python]
@dataclass
class SparseTensorChirho:
    indices_chirho: List[Tuple[int, ...]]  # Nonzero coordinates
    shape_chirho: Tuple[int, ...]          # Dimension sizes
\end{lstlisting}

Operations:
\begin{itemize}
  \item \textbf{AND} (element-wise): intersection of coordinate sets
  \item \textbf{OR} (element-wise): union of coordinate sets
  \item \textbf{Contraction}: hash-join on shared indices, OR-aggregate
\end{itemize}

\subsection{Boolean Contraction Algorithm}

\begin{algorithm}[ht]
\caption{Boolean Tensor Contraction}
\label{alg:contract_chirho}
\begin{algorithmic}[1]
\Require Tensors $T_A$, $T_B$ with shared indices $S$
\State Build hash table: $H[s] \gets \{a : (a, s) \in T_A\}$ for $s \in \text{proj}_S(T_A)$
\State result $\gets \{\}$
\For{$(b, s) \in T_B$ where $s \in H$}
  \For{$a \in H[s]$}
    \State result.add($(a, b)$)  \Comment{OR = set union}
  \EndFor
\EndFor
\State \Return result
\end{algorithmic}
\end{algorithm}

This is essentially a hash-join from databases, with OR (union) as the aggregation.

\subsection{Tensor Network Representation}

\begin{lstlisting}[language=Python]
@dataclass
class TensorNetworkChirho:
    tensors_chirho: List[TensorChirho]
    index_sizes_chirho: Dict[str, int]  # Variable -> domain size

    def contract_chirho(self, order_chirho: List[str]) -> TensorChirho:
        """Contract network following given variable elimination order."""
        # ... implementation details
\end{lstlisting}

% ============================================================================
\section{Evaluation}
\label{sec:evaluation_chirho}
% ============================================================================

\subsection{Heuristic Comparison}

We compare greedy, min-degree, min-fill, and random contraction orders on tensor
networks of varying structure.

\begin{table}[ht]
\centering
\begin{tabular}{lrrrr}
\toprule
\textbf{Network} & \textbf{Greedy} & \textbf{Min-Degree} & \textbf{Min-Fill} & \textbf{Random} \\
\midrule
Linear chain (3 tensors) & 20 & 20 & 20 & 25 \\
Star topology (4 tensors) & 45 & 45 & 45 & 68 \\
Complex (10 tensors) & 189 & 176 & 168 & 312 \\
\bottomrule
\end{tabular}
\caption{Estimated contraction cost (sum of intermediate sizes) by heuristic.}
\label{tab:heuristics_chirho}
\end{table}

\textbf{Finding:} Min-fill consistently produces the best orders, but greedy
is competitive (within 15\%) and faster to compute.

\subsection{Transitive Closure: Tensor vs.\ Datalog}

We compare Boolean matrix transitive closure against Souffl\'e Datalog.

\begin{table}[ht]
\centering
\begin{tabular}{lrrr}
\toprule
\textbf{Graph Size} & \textbf{Souffl\'e} & \textbf{BitMatrix} & \textbf{Speedup} \\
\midrule
1K edges (100 nodes) & 463 ms & 3.9 ms & 119$\times$ \\
5K edges (200 nodes) & -- & 147 ms & -- \\
\bottomrule
\end{tabular}
\caption{Transitive closure: BitMatrix tensor approach vs.\ Souffl\'e Datalog.}
\label{tab:datalog_chirho}
\end{table}

\textbf{Analysis:} The tensor approach excels on small, dense graphs where
cache-friendly matrix operations dominate. Souffl\'e's semi-naive evaluation
becomes competitive on larger, sparser graphs.

\subsection{Relational Append}

\begin{table}[ht]
\centering
\begin{tabular}{lrrr}
\toprule
\textbf{Operation} & \textbf{Tensor (Rust)} & \textbf{OCanren} & \textbf{Speedup} \\
\midrule
appendo (forward) & 1.2 $\mu$s & 45 $\mu$s & 37$\times$ \\
appendo (backward) & 8.5 $\mu$s & 320 $\mu$s & 37$\times$ \\
Composed appendo & 15 $\mu$s & 890 $\mu$s & 59$\times$ \\
\bottomrule
\end{tabular}
\caption{Relational append: tensor contraction vs.\ OCanren stream search.}
\label{tab:appendo_chirho}
\end{table}

% ============================================================================
\section{Discussion}
\label{sec:discussion_chirho}
% ============================================================================

\subsection{When Tensors Win}

The tensor approach excels when:
\begin{itemize}
  \item Domains are finite and reasonably small ($\leq 10^4$ values)
  \item Relations have moderate sparsity (1--10\% density)
  \item Multiple joins compose (contraction order optimization pays off)
  \item Hardware acceleration is available (GPU, FPGA)
\end{itemize}

\subsection{When Backtracking Wins}

Traditional backtracking excels when:
\begin{itemize}
  \item Domains are infinite or very large
  \item Pruning is effective (early failure detection)
  \item Memory is limited (backtracking uses $O(\text{depth})$ space)
  \item Solutions are sparse relative to search space
\end{itemize}

\subsection{Connection to Quantum Simulation}

The tensor contraction problem we solve is \emph{identical} to quantum circuit simulation.
A quantum circuit on $n$ qubits is a tensor network with $2^n$-dimensional indices.
Our contraction-order heuristics apply directly to quantum simulation, and vice versa.

This connection suggests that advances in quantum simulation algorithms
(e.g., from Google, IBM) could improve relational query processing.

% ============================================================================
\section{Related Work}
\label{sec:related_chirho}
% ============================================================================

\textbf{Tensor networks in AI.}
Tensor networks originate in quantum physics~\cite{tensor_networks_intro_chirho} and
have found applications in machine learning (tensor train decomposition~\cite{tensor_train_chirho})
and probabilistic inference.

\textbf{Join ordering in databases.}
The contraction order problem is equivalent to join ordering, studied extensively
in database query optimization. Hypertree decompositions~\cite{gottlob_hypertree_chirho}
provide theoretical bounds; practical systems use dynamic programming or greedy heuristics.

\textbf{Datalog engines.}
Souffl\'e~\cite{souffle_chirho} compiles Datalog to efficient C++ with semi-naive evaluation.
Scallop~\cite{scallop_chirho} extends Datalog with differentiable provenance semirings.
Our tensor approach is complementary: better for small, dense problems; Datalog
excels at large-scale program analysis.

\textbf{Constraint logic programming.}
CLP(FD)~\cite{clpfd_chirho} pioneered finite-domain constraint propagation in Prolog.
The tensor representation can be viewed as a batched, vectorized form of arc consistency.

% ============================================================================
\section{Conclusion}
\label{sec:conclusion_chirho}
% ============================================================================

We have shown that miniKanren relations correspond to sparse Boolean tensors and
relational composition corresponds to tensor contraction. This mapping:
\begin{itemize}
  \item Explains \texttt{conde} as tensor stacking (disjunction = union)
  \item Enables optimization via contraction-order heuristics
  \item Opens a path to hardware acceleration (GPU/FPGA tensor cores)
\end{itemize}

The contraction-order problem connects relational programming to quantum simulation
and database query optimization---three fields with shared algorithmic foundations.

\paragraph{Code Availability.}
Reference implementations:
\begin{itemize}
  \item \texttt{tensor\_network\_chirho.py} --- Boolean tensor composition
  \item \texttt{contraction\_order\_chirho.py} --- Heuristics (greedy, min-degree, min-fill)
\end{itemize}
Available at: \url{https://github.com/loveJesus/minikanren_1bit_chirho}

\paragraph{Companion Papers.}
This paper is part of a suite:
\begin{itemize}
  \item \paperA{} --- The core AND insight
  \item Paper B (this paper): Relations as Tensors
  \item \paperC{} --- Infinite domains via hash consing
  \item \paperD{} --- Tabling for termination
  \item \paperE{} --- FPGA acceleration
  \item \paperF{} --- Differentiable relaxation
\end{itemize}

% ============================================================================
\section*{Acknowledgments}
% ============================================================================

\standardacknowledgments

% ============================================================================
% Bibliography
% ============================================================================

\bibliographystyle{plain}
\bibliography{paper_b_chirho}

\vspace{1em}
\begin{center}
\textit{Soli Deo Gloria} $\chi\rho$
\end{center}

\end{document}
